% \documentclass{article}
\documentclass[leqno]{article}
% \documentclass[border=1cm]{standalone}
\usepackage[T2A]{fontenc}        
\usepackage[utf8]{inputenc}    
\usepackage[russian,english]{babel}
\usepackage{graphicx}
\usepackage{float}
\usepackage[hidelinks]{hyperref}
\usepackage{fancyvrb}
\usepackage{array}
\usepackage{booktabs}
\usepackage{siunitx}
\usepackage{tabularx}
\usepackage{longtable}
\usepackage{threeparttable}
\usepackage{amsmath}
\usepackage{mathtools}
\usepackage{bm}
\usepackage{tikz}
\usetikzlibrary{math}
% \usepackage{natbib}
% \usepackage[style=numeric]{biblatex}

% \usepackage{unicode-math}
% \setmainfont{TeX Gyre Pagella}
% \setmathfont{TeX Gyre Pagella Math}

% \addbibresource{learnlatex.bib} % file of reference info

\title{Computer Skills for Scientific Writing\\Диаграммы и чертежи в виде кода}
\author{Абрамян Артём Арменович}
\date{Февраль 2026}

% \newcommand\Triangle[2]{
%   \draw #1 coordinate(a) -- ++(0:#2) coordinate(b);
%   \draw (a) -- ++(60:#2) coordinate(c);
%   \fill (a) -- (b) -- (c) -- cycle;
% }

\begin{document}
  \maketitle
  \pagebreak

  \textbf{Цель работы:} Изучить возможности создания диаграм и рисования в LaTeX, разобраться с набором инструментов TikZ.\smallskip

  \section{Введение}
  \label{sec:intro}
  
Существует множество различных наборов инструментов для создания графических объектов в документе TeX. 
Самый известный и, вероятно, наиболее широко используемый из этих наборов инструментов — это TikZ, 
рекурсивная аббревиатура от «TikZ ist kein Zeichenprogramm» (нем. «TikZ — это не программа для рисования»). 
Из названия набора инструментов видно, что создание фигуры в TikZ не будет работать путем рисования фигуры, 
как это обычно делается в Paint или подобных программах для рисования. 
Для создания фигуры с помощью TikZ вам потребуется запрограммировать ее, используя специальные команды TeX. 
Здесь мы рассмотрим основы TikZ.

Три полезных источника информации и вдохновения для создания фигур TikZ:
• Статья о TikZ на Wikibooks: https://en.wikibooks.org/wiki/LaTeX/PGF/Ti
kZ.

• Статья о TikZ на CTAN с официальным руководством и другой документацией:
https://www.ctan.org/pkg/pgf.

• База данных TeXample TikZ, полная примеров и вдохновения: https://texa
mple.net/tikz/.

  
  \section{Основные компоненты TikZ}
  \label{sec:essential}
  
Для создания рисунка в TikZ вам сначала нужно загрузить соответствующий пакет в ваш файл TeX,
 то есть использовать \\usepackage{tikz} в преамбуле. 
 Существует широкий спектр опций и специальных библиотек TikZ, 
 которые также можно использовать и которые необходимо загрузить отдельно, но они нам не понадобятся в этом курсе для начинающих.
  Вы можете ознакомиться с доступными опциями и попробовать их самостоятельно, используя три ссылки, указанные выше.
Рисунок TikZ обычно помещается в среду LaTeX \\begin{figure}, как и другие рисунки. 
  Это полезно, поскольку позволяет присвоить рисунку номер и подпись, чтобы вы могли легко ссылаться на него. 
  Внутри среды рисунка вы можете начать создание рисунка TikZ, написав код \\begin{tikzpicture}, 
  после чего вы можете начать верстку рисунка. Как обычно,
вы закрываете среду кодом \\end{tikzpicture}.

В отличие от того, к чему вы привыкли при наборе текста и математических выражений в LaTeX, вам понадобится
один символ, который не соответствует напрямую тому, что вы можете увидеть
в вашем документе. Каждая строка кода TikZ должна быть закрыта точкой с запятой.
Давайте начнем с пошагового набора графика в TikZ. В конце этого
урока мы сможем набрать следующий рисунок.


    
  \section{Рисование линий}
  \label{sec:lines}
  
  График состоит из нескольких точек и прямых. Построение графика начинается с определения координат точек. 
  В TikZ это можно сделать двумя разными способами.
Декартова координата задается двумя числами (координатами x и y), разделенными запятой и заключенными в круглые скобки,
 например (-2,0) или (1,5,1,5). В TikZ координаты также можно задавать в полярной форме. 
 Это также работает, если задать два числа внутри круглых скобок (угол в градусах и длина), 
 разделенные двоеточием, например (180:2) или (45:2,1) для (приблизительно) тех же двух координат.
Примечание: если единица измерения не указана, по умолчанию используются сантиметры, координаты также можно задавать,
 например, (1 см, 2 пт).
Теперь мы можем нарисовать прямые линии между этими координатами, используя команду \\draw.


  \begin{tikzpicture}
\draw (-1,0) -- (3,10pt) -- (35:3);
\end{tikzpicture}

  

Помимо рисования прямых линий, мы также можем определять другие типы линий,
и этим линиям можно присвоить цвет, толщину или узор. Например,
используйте -| или |- вместо -- для рисования линий, которые сначала идут горизонтально, а затем вертикально,
или наоборот. Стиль линий можно изменить, задав параметры
команде \\draw. Это делается путем помещения этих параметров в квадратные скобки
после команды. В следующем примере мы рисуем стрелку вместо линии,
и задаем цвет другой линии.



\begin{tikzpicture}
\draw[->] (-1,0) -| (3,10pt);
\draw[red] (3,10pt) -- (35:3);
\end{tikzpicture}


Используйте исходный код TikZ, чтобы увидеть, какие типы опций возможны.
Мы можем рисовать не только прямые линии, но и кривые. 
Это можно сделать несколькими разными способами. В следующем примере мы представляем три варианта.


\begin{tikzpicture}
\draw (-1,0) to (5,1);
\draw[green] (-1,0) to[out=90,in=135] (5,1);
\draw[cyan] (-1,0) .. controls (0,-2) .. (5,1);
\end{tikzpicture}


Первое, что вы можете заметить, это использование `to` вместо `--` для соединения координат. 
Эта команда работает точно так же, но при использовании `to` мы можем добавить параметры к команде соединения.
 Для зеленой линии мы задаем два параметра: во-первых, угол (в градусах), под которым линия отклоняется от первой координаты,
  и во-вторых, угол, под которым линия достигает второй координаты. 
  Синяя линия — это так называемая кривая Безье, линия, определяемая ее начальной и конечной точками,
   а также рядом опорных точек (в данном случае одной), которые фактически притягивают линию к себе. 
   В TikZ такие линии могут быть определены с одной или двумя опорными точками. 
   В приведенном ниже примере мы приводим код для построения кривой Безье с двумя опорными точками.


\begin{tikzpicture}
\draw[dotted,gray] (-1,0) -- (5,1);
\draw (-1,0) .. controls (0,-2) and (4,2) .. (5,1);
\end{tikzpicture}




  \section{Nodes}
  \label{sec:node}

Следующий шаг на пути к верстке графика — это подписи, фрагменты текста и
числа, содержащиеся в рисунке. Для написания текста в рисунке TikZ нам нужен
объект, называемый узлом. Простейший пример узла — это текст, размещенный
на заданной координате, но узел также может представлять собой фрагмент текста, окруженный кругом,
прямоугольником или другой простой фигурой. Узел не обязательно должен содержать текст,
но в большинстве случаев текст является целью узла. Узел можно определить во время
рисования пути, что работает следующим образом:


\begin{tikzpicture}[scale=3]
\draw (0,0) node {hello} -- (1,1) node {world};
\end{tikzpicture}


Параметр [scale=3] в начале рисунка увеличивает масштаб рисунка в 3 раза. 
Командному узлу также можно задать параметры, например, разместить узел посередине линии, в конце линии и т. д.


\begin{tikzpicture}[scale=3]
\draw (0,0) -- (1,1) node[midway]{A} node[pos=0.75,above]{B} node[right]{C};
\end{tikzpicture}



Обратите внимание: при использовании команды `to` вместо `--` расположение узла `--` немного меняется. 
Приведенный ниже пример кода генерирует то же изображение, что и выше.


\begin{tikzpicture}[scale=3]
\draw (0,0) to node[midway]{A} node[pos=0.75, above]{B} (1,1) node[right]{C};
\end{tikzpicture}


Узлам также можно приписывать математические формулы, и, как уже говорилось, их можно
ограничить простой геометрической фигурой.


\begin{tikzpicture}[scale=3]
\draw (0,0) node[circle, draw]{$\sum_{i=1}^{n}n^2$} -- (1,1)
node[rectangle,draw]{$\frac{1}{\sqrt{2}}$};
\end{tikzpicture}


Как вы могли заметить, линия, соединяющая узлы, по-прежнему начинается и заканчивается
точно в заданных координатах, что выглядит не очень красиво. Эту проблему можно решить,
предварительно определив узлы и присвоив им метки, чтобы мы могли нарисовать линию от
узла к узлу.

\begin{tikzpicture}[scale=3]
% define nodes
\node[circle,draw] (label1) at (0,0) {$\sum_{i=1}^{n}n^2$};
\node[rectangle,draw] (label2) at (1,1)
{$\frac{1}{\sqrt{2}}$};
% draw the line
\draw (label1) -- (label2);
\end{tikzpicture}


Теперь мы рассмотрели достаточно примеров, чтобы понять, как можно оформить график в TikZ.
 Следующий фрагмент кода используется для генерации графика, показанного в начале этого урока.


 \begin{tikzpicture}[scale=2]
  % Define the nodes
\node[circle, draw] at (0,0) (a) {A};
\node[rectangle, fill] at (3,0) (b) {};
\node at (3,0.4) (blabel) {B};
\node[rectangle,rounded corners, draw] at (5,2) (c) {C};
% Draw the paths
\draw[->, green] (a) -- (b) node[midway, below,black]{2};
\draw[<->, blue] (a) to[out=45, in=135] (b);
\draw[->,red] (b)--(c);
\draw[yellow,dotted,very thick] (b) |- (c);
\draw[<-,cyan] (b) -| (c);
\draw[thick,black] (a).. controls (1,5) .. (c) node[midway, above]{$\frac{1}{2}$};
\end{tikzpicture}



  \section{Построение кривых}
  \label{sec:curves}

В TikZ можно не только рисовать линии между координатами. 
Также имеется встроенный способ построения графиков для широкого спектра различных функций. 
Например, для функции $y = x^2$ на интервале $[-2, 2]$.


\begin{tikzpicture}
\draw [domain=-2:2] plot (\x, {pow(\x,2});
\end{tikzpicture}





\begin{tikzpicture}[scale=1.5]
% Draw the x and y axis, label the axes and the origin
\draw[gray, ->] (-2,0) -- (2,0) node[right]{$x$}
 node[pos=0.53, below]{$O$};
\draw[gray, ->] (0,-1) -- (0,1) node[above]{$y$};
\draw[fill,gray] (0,0) circle [radius=1pt];
% Plot the curve
\draw[blue, thick] [domain=-2:2, samples=150] plot (\x,
 {cos(pi*\x r)})
node[right]{$y = \cos(x)$};
% Note: the r in the argument of the cosine signifies that
 we enter x in radians
\end{tikzpicture}


  \section{Работа с циклами}
  \label{sec:loops}

В TikZ также можно использовать циклы for, применяя команду \\foreach ... in ....
Таким образом, довольно сложное изображение можно очень быстро набрать.

\begin{tikzpicture}[scale=0.75]
\foreach \x in {0,1,2,3}
\draw[red,thick] (0,\x) circle [radius=\x+1];
\end{tikzpicture}


% Эта идея очень полезна для построения итеративных фигур, таких как треугольник Серпинского, 
% хотя для этого требуется специальная библиотека TikZ, с помощью которой можно определить рекурсивную функцию. 
% Следующий код показывает первые шесть итераций треугольника Серпинского,
% нарисованных в TikZ. Код также содержит некоторые элементы в преамбуле
%  (перед \\begin{document}), поэтому мы приводим полный код LaTeX для PDF-файла, содержащего только эту фигуру.


% Define a equilateral triangle with lower left corner at coordinate #1 and
% with length of the sides #2


% \begin{tikzpicture}
% \tikzmath{
% % Define the recursive function sierpinski
% function sierpinski(x, y, s, d) {
% if (d == 0) then {
% % Draw a triangle lower left corner at (\x, \y), length \s
% { Triangle{(x,y)}{s}; };
% } else {
% % Rescale the length of the sides and choose correct coords
% % for the next triangles
% u1 = 0.25*s;
% u2 = u1*sqrt(3);
% u3 = 0.5*s;
% sierpinski(x,y,u3,d-1);
% sierpinski(x+u3,y,u3,d-1);
% sierpinski(x+u1,y+u2,u3,d-1);
% };
% };
% % Let the length of the sides of the base triangle be 4, and generate 6 figures
% \S = 4;
% for d in {0,...,5}{
% % To situate all plots nicely under and next to each other, define the coords
% % of the lower left corners preemptively
% x = (S+1)*mod(d,2);
% y = int(d/2) * (S+1);
% sierpinski(x,-y,S,d);
% };
% \end{tikzpicture}



  \section{Упражнения}
  \label{sec:ex}

  \subsection{Graph 1}
  \label{subsec:first}


  \begin{tikzpicture}[scale=2]
  \node[circle, draw, fill=green!60, minimum size=8mm] at (0,1.8) {A};

\node[circle, draw, double, minimum size=8mm] at (-1.6,0.8) {B};
\node[circle, draw, fill=green!60, minimum size=8mm] at (-1.2,-0.8) {C};
\node[circle, draw, double, minimum size=8mm] at (0,-1.6) {D};
\node[circle, draw, fill=green!60, minimum size=8mm] at (1.2,-0.8) {E};
\node[circle, draw, double, minimum size=8mm] at (1.6,0.8) {F};

% --- Edges ---
\draw[blue, dotted, thick] (-1.6,0.8) -- (1.6,0.8)
  node[midway, above] {6};

\draw[blue, dotted, thick] (-1.6,0.8) -- (0,-1.6)
  node[midway, left] {2};

\draw[blue, dotted, thick] (1.6,0.8) -- (0,-1.6)
  node[midway, right] {4};

\draw[cyan, thick]
  (-1.6,0.8) to[out=-120, in=120] (-1.2,-0.8);

\draw[cyan, thick]
  (-1.2,-0.8) to[out=-30, in=-150] (0,-1.6);

\draw[red, thick]
  (0,1.8) to[out=180, in=90] (-1.6,0.8);

\draw[red, thick]
  (0,1.8) to[out=-20, in=90] (1.2,-0.8);

\draw[red, thick]
  (0,1.8) to[out=-90, in=120] (0,-1.6)
  node[midway, right] {$\sqrt{2}$};

\end{tikzpicture}

  \section{Заключение}
  \label{sec:conclusion}
  
  \textbf{Выводы:} В данной работе мы изучили возможности создания диаграм и рисования в LaTeX, разобрались с набором инструментов TikZ.

\end{document}